% TOP_PROBLEM: {"problemId": "problem.square-peg-problem-toeplitz-s-inscribed-square-conjecture", "canonicalTitle": "Square Peg Problem (Toeplitz's Inscribed Square Conjecture)", "releaseRank": 489, "primaryDomain": "geometry_topology", "primaryDomainLabel": "Geometry & topology", "displayStatus": "open_with_solved_subcases", "releaseStatus": "open_with_solved_subcases"}
% !TeX program = lualatex
% Definition and sources only; this file is not an LLM solution attempt.
\documentclass[11pt]{article}
\usepackage[margin=25mm]{geometry}
\usepackage{fontspec}
\setmainfont{Latin Modern Roman}
\usepackage{amsmath,amssymb,mathtools}
\usepackage{unicode-math}
\setmathfont{Latin Modern Math}
\usepackage{xurl,hyperref}
\hypersetup{colorlinks=true,urlcolor=blue}
\setcounter{secnumdepth}{0}
\setlength{\parindent}{0pt}
\setlength{\parskip}{5pt}
\setlength{\emergencystretch}{3em}
\begin{document}
\section*{\texorpdfstring{Square Peg Problem (Toeplitz's Inscribed Square Conjecture)}{Square Peg Problem (Toeplitz's Inscribed Square Conjecture)}}\label{489}
\subsection{Definitions and mathematical statement}
A Jordan curve is the image $\Gamma=\gamma(S^1)\subset{\mathbb{R}}^2$ of a continuous injective map from the circle. Let $J$ be rotation through $\pi/2$. The square-peg assertion is
\[
 \exists x,v\in{\mathbb{R}}^2,\ v\ne0:\quad
 x,\ x+v,\ x+v+Jv,\ x+Jv\in\Gamma.
\]
These are four distinct vertices of a Euclidean square. The curve need not be differentiable, rectifiable, or convex, and the square need not lie inside its bounded complementary region. Only the vertices are required to lie on the curve. A limiting configuration whose side length tends to zero does not supply a nondegenerate square.
\par\smallskip\textit{Formulation and concept references:} \href{https://arxiv.org/pdf/2203.02613}{[S1]}.

\subsection{Short English statement}
Does every continuous simple closed plane curve contain the four vertices of a nondegenerate square?
\subsection{Sources}
\begin{itemize}
\item[S1]Gregory R. Chambers, "On the Square Peg Problem," arXiv:2203.02613 [math.GT], March 2022.
\url{https://arxiv.org/pdf/2203.02613}.
\textit{Formulation source.}
\end{itemize}
\end{document}
